\subsection{Future work: the experiment that would test this directly}
\label{sec:future}
No public dataset images trained or tolerised human monocytes, so whether innate immune
memory has an imaging phenotype remains untested. The study that would test it
(Figure~\ref{fig:design}) uses primary human CD14$^+$ monocytes from independent donors in
five arms: medium, $\beta$-glucan (training), LPS (tolerance), and IFN-$\gamma$ and IL-4 as
polarisation anchors with known morphological and metabolic effects. Cells are exposed for
24\,h, washed and rested for 5 days, then restimulated with LPS on day 6. Label-free
brightfield time-lapse is recorded during the rest and for 0--24\,h after restimulation, at a
frame interval chosen to cover the horizons our forecasting analysis spanned (8\,s to
3.6\,min); NAD(P)H/FAD autofluorescence, with lifetime imaging where available, is acquired
on days 0, 1, 3 and 6. Ground truth is TNF and IL-6 protein at 4\,h and 24\,h after
restimulation, bulk RNA-seq at 4\,h, ATAC-seq on day 6 and lactate in the supernatant.

The analysis is fixed in advance: frozen video features (TiJepa, DINOv2 and V-JEPA~2) of
the rest-period movies, donor-held-out classification of naive, trained and tolerised wells,
and donor-held-out prediction of 24\,h TNF protein from pre-restimulation video, with
image-free baselines (well-level intensity and cell counts) as the reference; the primary
endpoint is the gain over that reference. The hypothesis is falsified if no frozen feature set
exceeds the reference with a donor-level confidence interval excluding zero. Effects
between donors are consistent in the transcript data (the SD of the trained-minus-untrained
TNF difference after restimulation is \val{ds.tnf.sd212} and \val{ds.tnf.sd141} log2 units
in the two datasets), but nothing constrains the size of an imaging effect. Assuming a
standardised donor-level effect of 1.0 (0.8, 1.5), a paired two-sided test at
$\alpha = 0.05$ needs \val{ds.n.d1.0} (\val{ds.n.d0.8}, \val{ds.n.d1.5}) donors for 80\%
power. Images would be deposited in the BioImage Archive under CC~BY, sequencing in GEO,
and cytokine and plate metadata alongside.
